\section{Introduction}
\label{sec:introduction}

\subsection{The unresolved attachment boundary}

End-to-end encryption changes the location at which an attachment can be
examined, but it does not remove the need for an examination decision. A relay
that follows the confidentiality model of a protected conversation does not
receive attachment plaintext. The endpoint, by contrast, eventually decrypts
the attachment and presents some representation of it to an operating-system
service, a codec, a preview generator, or an application renderer. The security
of the transport and the admissibility of the decrypted representation are
therefore distinct properties. Secure-messaging analyses establish valuable
communication guarantees, including confidentiality, authentication, and
forms of post-compromise security~\cite{unger2015,cohngordon2020}. They do not
establish that an arbitrary decrypted file has one unambiguous interpretation
or that every parser exposed by the receiving endpoint will process it safely.

This distinction matters because a file is not merely an inert sequence of
bytes. It is an input to a collection of recognizers whose views can differ.
The sender application may derive a type from a filename. A transport wrapper
may carry a declared media type. A preview service may inspect a prefix. A
codec may accept a valid initial object and ignore trailing bytes. A security
scanner may search for a set of signatures. The final application may invoke a
different parser with different stopping conditions. An adversary can exploit
the disagreement without defeating encryption. Encryption preserves the
attacker-selected bytes faithfully until they reach the endpoint that has the
keys and the parser attack surface.

Gateway scanning offers no general solution to this placement problem. A
gateway that sees plaintext changes the confidentiality model. A gateway that
sees only ciphertext cannot make a content-derived decision. Client-side
scanning can inspect plaintext, but a scanner verdict alone does not define the
representation later consumed by the renderer. Even a correct negative
signature result cannot establish that the scanner and renderer assigned the
same grammar to the bytes. Research on parser differentials, polyglot files,
and language-theoretic security shows that discrepancies among file-processing
components form an attack surface in their own right~\cite{jana2012,momot2016,
nail2014,everparse2019,ipg2023,koch2025,you2025}.

The term \emph{client-side} also requires a trust-purpose distinction.
Abelson et al. analyze client-side scanning in which an endpoint searches
cleartext for centrally designated targets and may report a match or its source
to an outside party~\cite{abelson2024}. The method studied here is instead a
user-benefiting endpoint defense. It reconstructs an attachment to reduce the
local parser exposure and returns no content-derived match report to the relay,
an authority, or another central service. The two designs share an execution
location, but not a beneficiary, reporting path, or surveillance capability.

The problem addressed in this manuscript is consequently narrower than
universal malware detection and stronger than isolated format sniffing. Given
an attacker-controlled attachment, the endpoint must decide whether it can
derive a new representation that belongs to a declared format profile, meets a
policy-selected reconstruction strength, survives independent complete
parsing, and can be published without exposing any other result. The relay must
remain blind throughout this procedure. The same decision must apply to an
outgoing attachment before encryption and to an incoming attachment after
decryption, because both directions can introduce hostile content into a
protected conversation.

\subsection{Why malware labels are not enough}

Terms such as virus, worm, Trojan, downloader, dropper, stager, ransomware, and
spyware describe different aspects of malicious software. Some terms concern
replication. Some concern deception or delivery. Others concern an operational
role or an objective after code execution. These labels do not by themselves
specify a byte-level property that a media reconstruction system can test. A
Trojan may be a native executable presented as a photograph. A downloader,
dropper, or stager may use a script hidden behind a misleading suffix.
Ransomware may arrive through a document, an archive, or an exploit-bearing
media sample. Spyware may be delivered as code, while a benign media file may
still reveal location or author metadata.
Conversely, a valid image can visually encode a command, a phishing message, or
arbitrary information without violating its image grammar.

The manuscript therefore uses malware labels and operational roles to organize
threats and corpus strata, not as outputs of \AFD. Each scenario is reduced to
an observable carrier property such as an executable prefix, trailing bytes, a
misleading portable name, a conflicting media type, an active embedded object,
excluded metadata, a malformed container, or a resource-exhaustion
construction. A successful result permits a claim about that property and the
configured policy. It does not permit the conclusion that the system has
identified a malware family or proved that retained pixels, samples, or codec
payloads contain no harmful information.

This claim discipline is necessary for two reasons. First, it prevents a
carrier-oriented method from being evaluated with labels that describe
post-execution behavior. Second, it makes each security statement falsifiable.
For example, a claim of complete output consumption is refuted by one accepted
output with unconsumed trailing bytes. A claim of metadata exclusion is refuted
by one accepted output containing a forbidden structural element. A broad
claim that a file is free of malware has no comparably precise test when the
system does not interpret all possible information encoded in the retained
media semantics.

\subsection{Research gap}

Existing work provides several pieces of the required foundation, but the
composition needed at an encrypted-messaging endpoint remains underspecified.
Parser-generation systems can connect grammars, parsers, and serializers.
Polyglot studies characterize cross-format ambiguity. Content disarm and
reconstruction systems replace or remove selected file components. Malware
engines detect known or learned indicators. Sandboxes constrain execution.
Secure-messaging protocols protect communications in transit. None of these
components alone answers the publication question posed at the endpoint:
\emph{which exact bytes may become available to the application after an
attachment has crossed an end-to-end encrypted channel?}

Three gaps are especially consequential. The first is an evidence gap. A file
extension, a declared media type, a magic-byte probe, and a complete parser do
not provide interchangeable evidence. Treating any one of them as decisive
creates a classifier-confusion path. The second is a reconstruction gap.
Container remultiplexing and semantic decode-and-reencode have different
security meanings. A structural path may remove metadata and reject trailing
objects while retaining the original compressed payload. A policy that requires
semantic reconstruction must not silently accept the weaker path. The third is
a publication gap. Failure must not leave partially reconstructed bytes,
unchecked temporary files, or a warning-labelled artifact that a downstream
caller can still render.

\AFD\ addresses these gaps as one endpoint-local method. It treats input
classification as bounded routing evidence, not acceptance. It records the
reconstruction level actually performed. It independently parses the output to
its exact terminal offset. It applies policy again to the reconstructed object.
It exposes deliverable bytes only when the final verdict is \Clean. These
conditions are enforced symmetrically around encryption and decryption, while
the relay continues to transport ciphertext.

\subsection{Objective and contributions}

The objective of this study is to define, implement, and pre-specify the
evaluation of an endpoint-local method that reduces attachment ambiguity while
preserving the confidentiality role of a blind encrypted relay. The object of
study is the attachment-processing boundary of an end-to-end encrypted
messaging system. The subject of study is the set of methods by which
format-profile convergence, canonical reconstruction, independent validation,
resource policy, and transactional publication can be composed into a
falsifiable delivery rule.

The residual novelty is the composition of five elements whose individual
mechanisms have prior art. The first is endpoint-local, user-benefiting
reconstruction at both plaintext boundaries of an end-to-end encrypted
messaging system, without a content-reporting path to the blind relay. The
second is convergence of exactly three routing sources, namely a normalized
portable name, a parsed caller-supplied media type, and bounded byte evidence,
followed by an independent complete output parser that authorizes delivery
rather than dispatch. The third is policy-indexed reconstruction strength, with
the performed level recorded and weaker reconstruction unable to satisfy a
stronger requirement. The fourth is independently validated clean-only
publication across in-memory, filesystem, and foreign-function interfaces. The
fifth is a pre-specified evidence contract that separates structural closure,
neutralization, false acceptance, false blocking, media fidelity, fault
behavior, and cost before empirical results are admitted.

The formal contribution is a separation between routing evidence and delivery
authorization. This separation yields an explicit acceptance predicate,
properties for complete consumption, canonical component closure, type
convergence, metadata exclusion, source resolution and provenance, resource
boundedness, independent candidate validation, and fail-closed result
construction. It also supports a policy order in which a stronger profile
cannot authorize an object solely because a weaker profile did. The
propositions presented later are analytical composition arguments under stated
assumptions. They are not represented as machine-checked proofs of codec
correctness.

The systems contribution is a Rust implementation whose architecture makes
the claim boundary inspectable. Handlers report the level actually performed.
A separate validation module parses rebuilt outputs. External codec processes
are invoked with time, output, diagnostic, protocol, and platform-dependent
resource constraints. Cancellation and panic paths terminate without a
deliverable object. Versioned policy and interface representations preserve the
decision across Rust, C-compatible, and application-facing boundaries.

The empirical contribution is the pre-specified evidence contract rather than
an early numerical claim. It defines benign canonical and noncanonical strata,
malformed inputs, classifier disagreements, true polyglots, boundary-smuggling
carriers, resource-exhaustion cases, unsupported active objects, and known
parser regressions. It separates format closure, neutralization, false
acceptance, false blocking, media fidelity, fault tolerance, and performance.
Numerical cells remain unavailable until a pinned server run passes the
declared release gates. This constraint prevents development-machine
observations from becoming publication results.

\subsection{Research question and scope}

The central research question is whether an endpoint-local reconstruction and
publication boundary can reduce measurable media-carrier ambiguity without
granting plaintext access or attachment authority to an intermediary. The
question decomposes into the more specific questions formalized in
\cref{sec:terminology-questions}. The study asks what evidence is sufficient
for dispatch, what additional evidence is necessary for acceptance, which
security properties follow from structural and semantic reconstruction, how
the policy behaves under strengthening, and how the implementation can be
evaluated without confusing parser evidence with behavioral malware detection.

The scope is intentionally restricted to configured media profiles. PDF,
office-document, archive, executable, script, package, shortcut, and
macro-bearing formats are outside the supported media grammar and are blocked
by the evaluated media-only profile. Phishing meaning, harmful imagery,
steganographic information retained in pixels or samples, and exploitation of
a previously compromised endpoint remain outside the claim. Decoder isolation
reduces consequences under stated operating-system controls, but it does not
prove an invoked decoder free of vulnerabilities.

The remainder of the manuscript proceeds from concepts to evidence.
\Cref{sec:terminology-questions} fixes terminology, units of analysis, research
questions, and falsification conditions. \Cref{sec:related-work} synthesizes the
relevant literature.
\Cref{sec:system-context,sec:malware-taxonomy,sec:extended-adversary}
establish system placement, malware-carrier taxonomy, and the adversary.
\Cref{sec:evidence-convergence,sec:formal-method,sec:policy-algebra,sec:publication-semantics}
define the method.
\Cref{sec:format-methods,sec:polyglots} instantiate it for supported formats
and cross-format constructions.
\Cref{sec:process-isolation,sec:implementation} trace the implementation.
\Cref{sec:corpus,sec:evaluation,sec:threat-campaigns,sec:results-framework}
pre-specify the empirical study. The discussion then distinguishes what the
construction establishes, what the future server evidence may establish, and
what remains outside both.
